\BookPart{Formal models of molecular computation}{formal}{%
Experiments show what molecules can do once; formal models show what they can
do in general. Part~IV builds the theory of molecular computation from small,
traceable examples: splicing and insertion systems inspired by recombination,
finite automata that recognise patterns, Watson--Crick automata that read both
strands of a duplex, and the constructions that make molecular systems as
powerful as any computer. Each model is paired with the physical process that
inspired it, so that a reader can see both what the abstraction captures and
what a real laboratory would still have to supply.

\textbf{Chapters in this part}
\begin{itemize}\setlength{\itemsep}{2pt}
\item \textbf{15\quad Splicing and insertion--deletion systems.} What languages can cut-and-paste rules generate?
\item \textbf{16\quad Languages, automata and molecular recognition.} How does a finite-state machine recognise a pattern, and how can molecules implement one?
\item \textbf{17\quad Watson--Crick and biochemical automata.} What can a machine reading both strands of a duplex compute, and how did enzyme-driven automata realise one?
\item \textbf{18\quad Universality and complexity models.} What does it take for a molecular system to compute anything computable, and at what cost?
\end{itemize}}
